% =========================================================
% config/commands.tex
% Custom command bersama seluruh buku.
% Jangan mendefinisikan command yang sudah ada di sini
% (lihat Aturan 2 pada Arsitektur_Buku_Ajar.md).
% =========================================================

% ---------------------------------------------------------
% Penanda istilah teknis
% ---------------------------------------------------------
\newcommand{\code}[1]{\texttt{#1}}          % kode / sintaks umum
\newcommand{\asm}[1]{\texttt{#1}}           % mnemonic instruksi assembly
\newcommand{\reg}[1]{\texttt{#1}}           % nama register
\newcommand{\flag}[1]{\texttt{#1}}          % nama flag
\newcommand{\dir}[1]{\texttt{#1}}           % direktif assembler
\newcommand{\file}[1]{\texttt{#1}}          % nama berkas
\newcommand{\eng}[1]{\textit{#1}}           % istilah asing

% ---------------------------------------------------------
% Bilangan
% ---------------------------------------------------------
\newcommand{\hex}[1]{\texttt{#1h}}          % heksadesimal: \hex{1234} -> 1234h
\newcommand{\bin}[1]{\texttt{#1b}}          % biner
\newcommand{\dec}[1]{\texttt{#1}}           % desimal

% ---------------------------------------------------------
% Penekanan
% ---------------------------------------------------------
\newcommand{\important}[1]{\textbf{Penting:} #1}
\newcommand{\note}[1]{\textbf{Catatan:} #1}
\newcommand{\term}[1]{\textbf{#1}}          % istilah yang didefinisikan

% ---------------------------------------------------------
% Singkatan OBE
% ---------------------------------------------------------
\newcommand{\CPL}[1]{\textbf{CPL#1}}
\newcommand{\CPMK}[1]{\textbf{CPMK#1}}
\newcommand{\SubCPMK}[1]{\textbf{Sub-CPMK #1}}

% ---------------------------------------------------------
% Tanda checklist
% ---------------------------------------------------------
\newcommand{\checkbox}{\ding{113}\hspace{0.35em}}
\newcommand{\checked}{\ding{51}\hspace{0.35em}}
\newcommand{\crossed}{\ding{55}\hspace{0.35em}}

% ---------------------------------------------------------
% Penanda revisi / TODO (jangan dipakai pada naskah final)
% ---------------------------------------------------------
\newcommand{\todo}[1]{\textcolor{red}{\textbf{[TODO: #1]}}}

% ---------------------------------------------------------
% Arah panah pada diagram
% ---------------------------------------------------------
\newcommand{\flowarrow}{\ensuremath{\rightarrow}}
